\paragraph{0D model.} The \texttt{ChamberCylinder} block of svZeroDSolver: a thick-walled
cylinder (inner radius $R_i = 19.0$~mm, outer radius $R_e = 33.0$~mm, length
$L_0 = 57.1$~mm) whose state is uniform along the axis and in $\theta$,
$r = R + \rho(R)$, $\theta = \Theta + \varphi(R) + \beta Z$, $z = Z + \eta(R)$,
discretized through the wall with \texttt{num\_elements} (ne) elements. All cases but the
free-length ones (Section~\ref{sec:closed}) use fixed length ($z = Z$ at the ends). The wall law is svMultiPhysics' \texttt{HolzapfelOgden} law
(the 0D block adopts the 3D law, including the ST91 volumetric penalty, here
$\kappa = 100\,\mu$), and the active case uses a prescribed fiber tension
$T(t) = 80~\mathrm{kPa}\,\sin^2(\pi(t-0.1~\mathrm{s})/0.4~\mathrm{s})$, $0.1 \le t \le 0.5$~s,
built exactly as svMultiPhysics' \texttt{UniformUnsteady} active stress.

\paragraph{3D model.} svMultiPhysics \texttt{struct} with quadratic tetrahedra (TET10,
``struct-p2''), the same material law, fiber directions given at the nodes, and a
follower pressure on the endocardium. Passive cases ramp the cavity pressure from 0 to
120~mmHg in 24 quasi-static steps of 50~s (5~mmHg). The mesh size is set by the number
of elements through the wall, nr ($h = (R_e - R_i)/\mathrm{nr}$).

\paragraph{Boundary conditions.} With circumferential or axial fibers the uniform 0D
state has the mirror symmetries of the cylinder, and the 3D model is an eighth: a
$90^\circ$ sector with symmetry planes over half the length, with $u_z = 0$ on the
mid-length plane and the end. Helical fibers break these symmetries, so the 3D model is
the full cylinder over the full length, of a tube twice as long ($2L_0$, both models;
Section~\ref{sec:helical}). As in the 0D model, both ends keep $u_z = 0$
and are otherwise free to twist and expand. In-plane dashpots on both end faces
(Robin, zero stiffness, damping $10^{-3}\,\mu/R_e\,\Delta t$) hold the three rigid-body
modes; at rest they carry no force. No 0D data enters the 3D model.

\paragraph{Comparison.} The 0D reference is the Richardson extrapolation of the 0D
block at ne and 2ne radial elements, compared at every load step (passive) or time step
(active). Kinematic quantities: cavity volume $V$, the through-wall profiles
$\rho(R)$, $\varphi(R)$, $\eta(R)$ at the radial node levels of the 3D mesh (averaged
over $\theta$ and $Z$), the wall thickness change $h$, and the twist per unit length
$\beta$. Struct fields are computed from the 3D displacement (Green--Lagrange strains,
$J$, fiber stretch, Cauchy stresses). With helical fibers the 3D tube differs from the
uniform 0D state near its free ends, so only the mid-length section is compared: the
middle element layer, $\beta$ as the local twist rate, and $V = \pi(R_i + \rho(R_i))^2 L$
as in the 0D model.

\paragraph{Active contraction.} The 3D cylinder replaces the 0D block in the circulation
of \texttt{tests/cases/chamber\_cylinder.json} (pressure source, valves, RC vessels,
1~s period), coupled through svMultiPhysics' \texttt{svZeroDSolver\_interface}
(implicit coupling, the cavity pressure as a coupled Neumann condition). The first
cycle (from 0.02~s on) is compared with the 0D block in the same circulation.

\paragraph{Computing.} Runs used svMultiPhysics in the \texttt{simvascular/solver} Docker
image (pinned digest), locally (emulated x86 on Apple Silicon) and on Yale's Bouchet
cluster (the same image in Apptainer, native, about five times faster per core;
results bitwise identical to the local runs for the same number of MPI ranks).
